\chapter{Literature Review}
\label{ch:literature}

This chapter situates APEX in the literature on online assessment, automatic grading of programming, and the security and usability concerns specific to classroom portals. The review is selective rather than exhaustive: the goal is to show what existing systems do well, what they trade away, and where APEX places itself.

\section{Online assessment in higher education}

The first generation of online assessment platforms grew out of the learning-management system tradition. Moodle~\cite{moodle2002}, originally released in 2002, integrated quiz delivery and gradebook into a single PHP application; Sakai followed in 2004 with a heavier Java stack; Canvas LMS~\cite{canvas} (2008) and Blackboard refined the form factor for large institutions. \Cref{fig:F2.1} draws a timeline of representative platforms.

\begin{figure}[H]
  \centering
  \includegraphics[width=0.95\linewidth]{F2.1_platform-timeline.pdf}
  \caption{Selected online-assessment platforms by release year.}
  \label{fig:F2.1}
\end{figure}

The strength of LMS-style platforms is their breadth: they cover the entire course lifecycle from announcement to gradebook. The weakness, for our purposes, is that programming-style assessment is bolted on. Moodle's quiz module supports MCQ and short-answer questions natively, but coding questions either require third-party plugins (CodeRunner~\cite{coderunner} is the best-known) or external tools through LTI. The result is a workflow where a teacher who wants to mix coding, MCQ, and written questions in a single timed exam typically settles for one format and supplements it elsewhere.

The pedagogical literature has long documented this gap. Douce, Livingstone and Orwell~\cite{douce2005autoassess} surveyed automatic test-based assessment systems in 2005 and observed that course management and code execution were typically handled by separate systems with weak integration. Ihantola et al.~\cite{ihantola2010autograding} returned to the same theme five years later and found the situation broadly unchanged: classroom features and grading features lived in different products, and the work of glueing them together fell on the instructor.

\section{Online judges and auto-grading}

The online-judge tradition emerged in parallel from competitive programming, where the performance and security of the sandbox dominated other concerns. UVa Online Judge, Codeforces, and HackerRank~\cite{hackerrank} all run student code against a battery of test cases inside an isolation layer. Wasik et al.~\cite{wasik2018survey} surveyed the field in 2018 and identified three recurring components: the contest/exam scaffold, the language-aware compile/run sandbox, and the verdict aggregator that turns per-test outcomes into a score. Caceres, Mateos, et al.~\cite{caceres2014online} surveyed the web frameworks behind such systems and noted that the durable winners typically separate the sandbox from the scaffold by a thin HTTP API.

Judge0~\cite{judge0,dosilovic2018judge0} is the most prominent open-source instance of that pattern: a Ruby + Rails service that wraps the \texttt{isolate} sandbox~\cite{isolate} and exposes a REST endpoint for "compile, run with stdin, return stdout, stderr, time, memory". \texttt{isolate} is the same sandbox used by the International Olympiad in Informatics, and provides the kernel-level isolation that running untrusted student code requires.

Judge0's appeal is precisely that it does one thing well. APEX builds on this --- the entire backend interaction with the sandbox lives in \texttt{internal/judge0/client.go} and consists of a single \texttt{POST /submissions?wait=true} call --- and frees APEX from having to ship the most security-sensitive component itself.

The remaining problem is the scaffold. Plain online judges typically lack class management, written-answer grading, and the per-class assignment workflow that a working classroom needs. Bridging that gap has been the subject of repeated academic projects (Web-CAT, Marmoset, AutoLab, OnlineQuiz, and a long tail of one-off institutional systems). APEX is closest in spirit to CodeRunner~\cite{coderunner} (which delivers programming questions through Moodle), Web-CAT (which adds a peer-review and TA-grading layer), and AutoLab (which targets large CS courses). Where APEX differs is the deliberate choice to integrate three answer formats (coding, MCQ, written) under a single attempt, with a single timer, a single autosave path, and a single notification.

\paragraph{Output diffing vs.\ unit tests.} Auto-grading approaches divide roughly into two families: input/output diffing (the student program reads stdin and writes stdout; the grader compares against an expected output) and unit-test execution (the student writes a function; the grader runs a test harness). Both have a long history in the literature~\cite{douce2005autoassess,ihantola2010autograding,wasik2018survey}. APEX uses I/O diffing for the same reasons most online judges do: it is language-agnostic, simple to specify, and resilient to small student deviations from a method signature. The price is that I/O diffing forgives some bugs (a method that prints the answer in a slightly different format) and rejects some correct solutions (a trailing newline mismatch). \texttt{judge0.normalizeOutput} is APEX's response to the second of those: line-end normalisation and trailing-whitespace trimming reduce false negatives without compromising detection of genuinely wrong outputs.

\section{Authentication and security}

JWT-based authentication~\cite{rfc7519} has become the default for stateless web APIs. The standard's strengths --- self-contained, signature-verifiable, no shared session store --- are paid for by the difficulty of revocation and by the well-known foot-guns: the "alg=none" forgery, weak HMAC secrets, and the temptation to put authorisation decisions inside the token rather than re-checking against the database. OWASP's Top 10~\cite{owasp-top10} catalogues these patterns under broken access control and security misconfiguration; CWE-639~\cite{cwe-639} (insecure direct object reference), CWE-79~\cite{cwe-79} (cross-site scripting), and CWE-352~\cite{cwe-352} (cross-site request forgery) cover the three vulnerabilities most relevant to a classroom portal.

OAuth 2.0~\cite{rfc6749} and OpenID Connect~\cite{oidc-core} are the standards behind the Google sign-in path. Google's Identity Services~\cite{gids} simplify the client-side dance but leave the verification of the ID token to the relying party --- in APEX's case to \texttt{idtoken.Validate} from Google's Go SDK. Storing only the resulting \texttt{google\_id}, \texttt{email}, and a verified flag avoids the complexity of refresh-token storage that full OAuth would require.

bcrypt~\cite{provos1999bcrypt} remains the recommended password hashing function in the absence of newer constructions (Argon2 is preferred where available; bcrypt is universally available across the Go ecosystem and is sufficient for the threat model). The constant-time comparison and bounded compute cost are the relevant properties.

\section{Cross-site scripting and content rendering}

Stored cross-site scripting~\cite{cwe-79} is the dominant client-side risk in any system where one user's input is rendered into another user's browser. Markdown-based renderers complicate the picture: a vanilla \texttt{react-markdown} pipeline is safe by default because it does not reparse raw HTML, but plugins like \texttt{rehype-raw} explicitly opt back in. The literature on safe Markdown rendering generally recommends either a strict subset (no inline HTML at all) or a sanitisation pass over the resulting HTML tree. APEX takes the strict-subset approach: \texttt{rehype-raw} is removed in commit \texttt{2eaac88}, and the resulting Markdown grammar is intentionally less expressive than the GitHub-flavoured Markdown it derives from.

\section{Usability of assessment systems}

Nielsen's heuristics~\cite{nielsen-heuristics} and Shneiderman's golden rules~\cite{shneiderman-rules} are the canonical heuristic frameworks for interface evaluation. They translate into a small set of concrete expectations for an exam-taking SPA: a visible system status (the countdown timer, the autosave indicator), error prevention (refusing to publish without a start time, refusing to submit twice), recognition over recall (the question rail with visited markers), and undo where it is safe (resume after autosave).

For quantitative usability, the System Usability Scale~\cite{brooke1996sus} is the default tool. SUS is a ten-item Likert questionnaire that yields a score between 0 and 100; the literature treats 68 as the industry average, with 80 considered "excellent". \Cref{ch:results} reports the results of a small SUS study against APEX.

\section{Comparable systems}

\Cref{tab:peer-comparison} compares APEX with Moodle (with CodeRunner), Canvas LMS (with the third-party PrairieLearn extension), HackerRank for Schools, and CodeRunner standalone. The criteria are the ones that the requirements in \Cref{ch:requirements} bring to the front. \Cref{fig:F1.2} drew the same data as a radar chart in the introduction.

\begin{table}[H]
  \centering
  \caption{Feature coverage at a glance (subjective; 0 = absent, 5 = first-class).}
  \label{tab:peer-comparison}
  \small
  \begin{tabularx}{\linewidth}{X c c c c c}
    \toprule
    Capability & APEX & Moodle+CR & Canvas & HackerRank & CodeRunner \\
    \midrule
    Live coding (Monaco editor)               & 5 & 2 & 2 & 5 & 4 \\
    Auto-grading via sandbox                  & 4 & 2 & 2 & 5 & 4 \\
    MCQ + written + coding in one exam        & 5 & 5 & 5 & 2 & 1 \\
    Manual grading queue                       & 5 & 5 & 5 & 1 & 1 \\
    Class invite codes + rosters              & 4 & 5 & 5 & 2 & 4 \\
    Reminders (in-app + email)                 & 4 & 5 & 5 & 3 & 3 \\
    Self-host (simple)                         & 5 & 3 & 1 & 1 & 4 \\
    \bottomrule
  \end{tabularx}
\end{table}

The numbers are subjective and meant to summarise the qualitative comparison rather than to claim measurement. The pattern is consistent: LMS-derived platforms (Moodle, Canvas) are stronger on classroom workflow but weaker on first-class coding; judge-derived platforms (HackerRank, CodeRunner) are the inverse. APEX's positioning is to cover both halves at "good enough", at the cost of being newer and less battle-tested than any of the comparators.

\section{Plagiarism detection and proctoring}

Plagiarism detection is largely orthogonal to the grading subsystem and is delivered by separate tools. MOSS~\cite{moss} (Stanford) and JPlag~\cite{prechelt2002jplag} are the two canonical open systems; both ingest a corpus of student source files and report similarity matches. Integration with a portal is typically a teacher-triggered batch action rather than an online check, and APEX does not yet ship that integration --- it is on the roadmap (\Cref{ch:conclusion}).

Proctoring --- webcam capture, screen sharing, environment locking --- has its own large literature and a parallel ethical debate. APEX deliberately does not pursue heavyweight proctoring; the most a future iteration is likely to add is non-invasive signals (tab-blur, fullscreen exit, paste detection) surfaced to the teacher rather than enforced.

\section{Where APEX places itself}

Three positioning claims summarise this chapter:

\begin{itemize}
  \item \emph{Integration over feature breadth.} APEX integrates the three exam formats more tightly than any LMS-bolted-on coding plugin and integrates classroom workflow more tightly than any pure online judge.
  \item \emph{Self-host before scale.} APEX's deployment story (\Cref{ch:deployment}) targets a single-instructor or single-department deployment. The price is the absence of features (per-tenant isolation, multi-region failover) that a commercial product would need.
  \item \emph{Defensive defaults over feature-flagged hardening.} The security review (\Cref{ch:security}) is part of the baseline rather than an enterprise upsell.
\end{itemize}

The next chapter --- \Cref{ch:requirements} --- formalises what this positioning means as a requirements list, and \Cref{ch:design} translates the requirements into the design that the rest of the book documents.
